\documentclass[11pt]{article}

\usepackage[T1]{fontenc}
\usepackage[margin=1in]{geometry}
\usepackage{amsmath}
\usepackage{amssymb}
\usepackage{booktabs}
\usepackage{array}
\usepackage{graphicx}
\usepackage{xcolor}
\usepackage{caption}
\usepackage{microtype}
\usepackage[colorlinks=true,linkcolor=blue!45!black,citecolor=blue!45!black,urlcolor=blue!45!black]{hyperref}

% Float packing: let tables share pages with more text instead of drifting.
\renewcommand{\topfraction}{0.9}
\renewcommand{\bottomfraction}{0.6}
\renewcommand{\textfraction}{0.07}
\renewcommand{\floatpagefraction}{0.8}
\setcounter{totalnumber}{4}
\setcounter{topnumber}{3}

\newcommand{\code}[1]{\texttt{#1}}
\newcommand{\usff}{\code{us54}}

\title{Inference From a Miss:\\ What Silence Licenses, and Two Corrections That Subtract}
\author{Allen Hsieh\\[2pt] {\small FishAI project}}
\date{August 2026}

\begin{document}
\maketitle

\begin{abstract}
\noindent
The owner of this project raised a specific worry about the engine's inference: if you
ask for a low spade, miss, and the target never asks you for a low spade back, do not
conclude they hold none. This paper answers that worry and then generalises it into the
full matrix for choosing an ask. The engine does not make the inference and structurally
cannot --- an absence is not a log event, so nothing in the ingestion \code{switch} can
key on one --- demonstrated by a treatment/control pair in which only the named card's
candidate set moves and the other four are bit-identical. The residual error runs
\emph{opposite} to the worry: in exactly the owner's position the engine believes
$14.0\%$ where the truth is $2.0\%$ ($N = 19{,}003$), too generous to the silent seat.
Calibration over $1{,}395{,}876$ scored asks from $300$ \usff{} games puts the overall
bias at $-0.0109$ --- optimistic, which is the signature of an engine that is \emph{not}
over-inferring --- with the certainty tier exact: $16{,}680$ asks called certain hit
$16{,}680$ times. Stratified, silence alone is worth $-0.0081$ over $1{,}176{,}652$ asks
and is genuinely near-uninformative, while silence \emph{after a miss} is a different
object --- such a seat holds a card of the set $6.3\%$ of the time against $66.3\%$ for a
silent seat never missed on. Two real gaps run the other way, both under-reading
\emph{positive} evidence: a published row-6 licence is under-priced by $8.4$ points
($0.2386$ believed against $0.3221$ observed, $N = 23{,}345$), because a deal-time
constraint is dropped the moment it is satisfied or exhausted --- exactly when the seat
has been shown to hold a card of the set; and the conceded-turn proxy is hand size, which
correlates $-0.147$ with the cards the conceded seat then takes against $+0.421$ for its
published reach. Conditioning on the licence at strength $\lambda = 0.60$ removes the
calibration bias almost exactly ($-0.0002$). The headline is what that bought. On one
$400$-pair common baseline, defusal alone is worth $+1.915 \pm 0.315$ sets per duplicate
pair, the calibration fix alone $+2.068 \pm 0.293$, and \textbf{both together
$+1.565 \pm 0.317$}: the stack ranks below either part, which is the shape of two
substitutes reached different ways rather than of complements. That ordering is
\emph{suggestive, not established}. Every interval quoted is an arm against the baseline;
none is on a difference between arms, and the smaller gap the reading rests on
($0.350$ sets) is under these arms' own $0.42$--$0.47$ detection floor. The arms do share
one bank and one baseline, so the comparison is paired and sharper than a between-bank one
would be --- but the paired interval was not computed, and at this $N$ the ordering is
reported as a signal to re-measure rather than as a settled result. Only a common-baseline
design raises the question at all --- measured each against its own baseline the
calibration fix reads as a flat $+0.04 \pm 0.41$. The same $400$-pair scale is where this instrument is weakest, and the
paper reports that in the body: the defusal quantity measures $+1.62$/$+1.66$ at
$4{,}000$ pairs, and the first explanation offered for the gap was tested and refuted.
A separate observation --- that two or three sets get ``ask-traded'' --- is confirmed at
$88.1\%$ against a $21.4\%$ availability-matched control and then shown not to be
reciprocation at all: $99.1\%$ of those replies are the hit-probability argmax. The
recommendation is to ship neither change.
\end{abstract}

\section{Introduction}
\label{sec:intro}

Every ask in Literature (Canadian Fish) is a question put in public, and every answer is
heard by the whole table. That makes the game a laboratory for a narrow question with a
wide reach: \emph{what does an event license you to conclude, and what does a
non-event?}

The question arrived here in a concrete form. The project's owner, watching the engine
play, wrote down a worry:

\begin{quote}
``If you have a low spade, you ask for a low spade but don't get it, and they don't ask
for a low spade back, don't automatically infer that they don't have any low spades.''
\end{quote}

It is a good worry. It names a real failure mode of enthusiastic deduction engines --- the
slide from ``no evidence that they hold one'' to ``evidence that they hold none'' --- and
it names it in the exact position where the slide is most tempting, because the miss
\emph{did} certify something and the silence sits right next to it.

This paper answers the worry, and then does the thing the worry was a special case of. It
audits every path by which the engine eliminates a candidate and asks of each whether it
is sound (Section~\ref{sec:audit}). It measures whether the engine over-infers anywhere,
by scoring its own stated hit probability against ground truth over $1{,}395{,}876$ asks
(Section~\ref{sec:calibration}). It measures what silence is actually worth, in the
owner's exact position and in three neighbouring ones (Section~\ref{sec:silence}). It
reports the two places where the engine's inference \emph{is} wrong --- both of them the
opposite of the feared error (Section~\ref{sec:gaps}) --- and one behaviour that looks
wrong from the outside and is not (Section~\ref{sec:clustering}). It lays out the full
ask matrix the owner asked for, twelve factors with their grounding, sign, and whether
they are priced today (Section~\ref{sec:matrix}). And it closes on the result that
changed what the project would do (Section~\ref{sec:substitution}): fixing the
calibration error is worth about two sets per duplicate pair, the shipped defusal
heuristic is worth about two sets per duplicate pair, and \emph{doing both measures lower
than doing either} --- an ordering the design makes visible and this bank size does not
resolve.

Three findings are worth stating before the machinery.

\textbf{The engine does not make the owner's mistake, and cannot.} Not because it is
careful, but because an absence is not representable in the layer that would have to act
on it (Section~\ref{sec:absence}). The measured residual runs the other way: the engine
is too generous to the seat that has gone quiet, not too harsh.

\textbf{Silence and silence-after-a-miss are different objects.} The first is worth
$0.8$ points of bias over $1.18$ million asks --- the owner's principle, validated. The
second is worth $60$ points of set-level probability. The compound event is informative
even though neither part is, and the engine under-uses it. It is deliberately not acted
on, for reasons Section~\ref{sec:notacting} gives in full.

\textbf{Two corrections for the same evidence should not be expected to add.} The defusal
term the project already ships \cite{concession} and the calibration fix this paper derives
are, at bottom, one reordering reached two ways --- an argument from mechanism, which the
measurement then agrees with: stacked, they measure lower than either alone. The ordering
is suggestive rather than resolved at this bank size, for the reasons
Section~\ref{sec:subresult} sets out. The design point stands regardless of how the
ordering settles, and Section~\ref{sec:lesson} states it once, in general.

\section{The rules the inference rests on}
\label{sec:rules}

Literature is played by six players in two teams of three (seats $0,2,4$ versus $1,3,5$)
\cite{pagat}. Under \usff{} \cite{rules54} the deck is 54 cards partitioned into nine
\emph{sets} (half-suits) of six; nine cards are dealt to each player; the game ends when a
team has been awarded five sets. The engine is pure, deterministic TypeScript, and one
seed reproduces a byte-identical game \cite{repo,v05paper}.

The inference in this paper uses a handful of rows of the \usff{} decision table
\cite{rules54}, cited by number throughout.

\begin{center}
\small
\begin{tabular}{@{}c p{0.84\textwidth}@{}}
\toprule
row & rule \\
\midrule
5  & \textbf{Ask target.} The target of an ask is always an opponent; a teammate may
     never be asked. \\
6  & \textbf{Ask licence.} An asker must hold at least one card of the asked set. \\
7  & \textbf{Own card.} You may not ask for a card you hold. \\
9  & \textbf{Hit.} The card transfers; the asker keeps the turn and asks again. \\
10 & \textbf{Miss.} The turn passes to the player who was asked. \\
12 & \textbf{Declare content.} Name the set and the exact holder of all six cards; every
     assignment must be a seat on the declarer's own team. \\
17 & \textbf{Information.} Card counts are public, with a full public log of asks,
     results, and declares. \\
\bottomrule
\end{tabular}
\end{center}

Two consequences do most of the work below. First, rows 6 and 7 make every ask a
\emph{broadcast}: it publishes permanently that the asker holds at least one card of the
set, and that the asker lacks the named card. Second, rows 9 and 10 make every ask a
\emph{bet with a tempo stake}: a hit takes a card and keeps the turn, a miss ends the turn
and hands it to a named opponent. The companion document treats the broadcast and the
concession \cite{concession}; this paper treats the bet, and what the broadcast lets a
listener conclude. The public channel both of them read from --- what a log of asks,
results and declares does and does not expose --- is specified in \cite{obspaper}.

\section{The direct answer: an absence is not an event}
\label{sec:absence}

\subsection{Why the engine cannot make the inference}

FishAI never treats silence as evidence. After $A$ asks $B$ for the $3\spadesuit$ and
misses, the engine concludes exactly one negative fact --- \emph{$B$ lacks the
$3\spadesuit$} --- and nothing whatever about $B$'s
$2\spadesuit/4\spadesuit/5\spadesuit/6\spadesuit/7\spadesuit$. $B$'s later asks into other
sets move those beliefs by zero \cite{asking}.

The reason is structural rather than careful, and the distinction matters for how much
confidence the result deserves. The ingestion routine in \code{knowledge.ts} is a
\code{switch} over events that \emph{are} in the log \cite{repo}. There is no scan for
``who has not asked into set $S$'', no silence counter, no candidate-mask decay. An
absence is not representable in that layer, so it cannot be acted on there. A guarantee of
that shape is stronger than a guarantee that rests on a careful predicate, because there
is no predicate to get wrong.

\subsection{Demonstrated, not asserted}
\label{sec:demo}

The claim is executed rather than argued: a treatment in which the $3\spadesuit$ ask
happens against a control in which it does not, with seat 1 then asking four more times
and never into low spades \cite{asking}.

\begin{table}[ht]
\centering
\caption{Candidate sets and per-seat probabilities after the treatment ask, against the
control that omits it \cite{asking}. Only the named card moves.}
\label{tab:demo}
\small
\begin{tabular}{@{}l l l r r@{}}
\toprule
card & candidates (treatment) & candidates (control) & $p(\text{seat }1)$ treat.
& $p(\text{seat }1)$ ctrl. \\
\midrule
$3\spadesuit$ & $[2,3,4,5]$   & $[1,2,3,4,5]$ & 0.0000 & 0.2000 \\
$5\spadesuit$ & $[1,2,3,4,5]$ & $[1,2,3,4,5]$ & \textbf{0.2000} & \textbf{0.2000} \\
$6\spadesuit$ & $[1,2,3,4,5]$ & $[1,2,3,4,5]$ & \textbf{0.2000} & \textbf{0.2000} \\
$7\spadesuit$ & $[1,2,3,4,5]$ & $[1,2,3,4,5]$ & \textbf{0.2000} & \textbf{0.2000} \\
\bottomrule
\end{tabular}
\end{table}

Table~\ref{tab:demo} is the whole answer to the literal question. The named card's
candidate set loses seat 1 --- which is row 7 and the logged miss, both sound --- and the
other four rows are bit-identical between the two runs. Four further asks by the silent
seat, into other sets, move nothing.

\subsection{The residual runs the other way}
\label{sec:residual}

The interesting part is what the engine believes in that position, measured against
ground truth. In the owner's exact case --- I missed on this seat, and it has never asked
back into the set --- the engine believes $14.0\%$ and the truth is $2.0\%$
($N = 19{,}003$) \cite{asking}.

It is \emph{too generous} to the silent seat, by twelve points, not too harsh. It never
writes that seat off; if anything it should write it off slightly more than it does. So
the owner's principle is right, the engine already honours it, and it honours it with
room to spare. Section~\ref{sec:silence} takes that twelve-point gap seriously as a
quantity, and Section~\ref{sec:notacting} explains why the project declines to spend it.

\section{Every way a candidate can be eliminated, and whether it is sound}
\label{sec:audit}

An engine that does not infer from absence can still infer unsoundly from presence. The
audit below walks every path by which a candidate seat is removed from a card's candidate
set, and asks of each what licenses it \cite{asking,repo}.

\begin{table}[ht]
\centering
\caption{The elimination-path audit \cite{asking}. Not one trigger is keyed on an
absence.}
\label{tab:audit}
\small
\begin{tabular}{@{}p{0.30\textwidth} p{0.47\textwidth} l@{}}
\toprule
trigger & justification & verdict \\
\midrule
a miss --- the target is cleared for the \textbf{named card} & the logged event
& \textbf{sound} \\
a miss --- the asker is cleared for the named card & row 7: you may not ask for a card you
hold (correctly disabled when \code{askOwnCardAllowed}) & \textbf{sound} \\
historical count exhaustion & pigeonhole on log-replayed counts; disabled when the log is
windowed & \textbf{sound} \\
current count exhaustion & pigeonhole on the public counts of row 17 & \textbf{sound} \\
own-hand injection & the viewer's own hand is ground truth & \textbf{sound} \\
singleton collapse, count forcing, constraint forcing, claim reveal & positive fixes,
downstream of the above or read from a claim's revealed holders (row 12) & \textbf{sound}
\\
\bottomrule
\end{tabular}
\end{table}

Three notes on Table~\ref{tab:audit}.

\textbf{The bounded arm inherits soundness rather than risking it.} The v1.5 bit-budget
fact pool \cite{v15paper} carries \code{lacks-card} and \code{no-basis} facts, which are
compressions of already-certified eliminations; and forgetting can only \emph{widen} a
candidate set, never narrow one. A bounded seat is therefore less certain than a
full-memory seat and never differently certain.

\textbf{The refined probability's only failure mode is over-confidence.}
\code{refinedHitProbability} short-circuits at zero and is monotone non-decreasing, so it
cannot manufacture a false negative --- it cannot be the mechanism that writes a seat off.
Section~\ref{sec:calibration} then measures whether it is over-confident in practice, and
finds it is not.

\textbf{Two places are sound-but-fragile, and neither is reachable at the shipped hard
tier.} A first-write-wins rule could persist a wrong certainty on \emph{inconsistent}
input, but is reachable only under the easy tier's six-event window, which also runs
without constraints; and \code{clearCand} refuses to empty a candidate set, which is the
anti-unsoundness guard itself, erring toward a superset. \textbf{No bug was found and no
code change is warranted} by this audit \cite{asking}.

\subsection{The one defect the audit did surface}
\label{sec:staleness}

The audit is about elimination, and elimination is sound. The defect it surfaced is about
\emph{retention}, and it is the cause of the largest gap in this paper.

\code{knowledge.ts} records an ask as a deal-time set constraint --- \emph{this seat holds
at least one card of this set} --- and \textbf{drops that constraint the moment it is
satisfied or exhausted} \cite{repo,concession}. Read plainly: the constraint is discarded
exactly when the seat has been \emph{shown} to hold a card of the set, which is precisely
when the fact is most informative. \code{refinedHitProbability} iterates those
constraints, so it goes blind where the evidence is strongest.

This is not a soundness bug. Nothing false is concluded; something true is forgotten. It
is nonetheless the origin of the $8.4$-point under-pricing of
Section~\ref{sec:underpricing}, and the threat model already routes around it: the header
of \code{threat.ts} names the same staleness and reads the licence off the public log
instead, which roughly doubles the share of harvest threats a seat can see
($6.0\% \to 11.7\%$) and takes detection of the canonical five-and-one position from
$7.1\%$ to $35.0\%$ \cite{concession,repo}. The fix had never reached the probability.

\section{Is the engine over-inferring anywhere? Calibration says no}
\label{sec:calibration}

The audit establishes that each elimination path is licensed. It does not establish that
the resulting probability is honest, and the two are different claims: an engine can
reason from sound facts to over-confident numbers. So every legal ask at every ask
decision was scored with the number \code{pickAsk} actually decides on, against ground
truth read from the hands --- $300$ \usff{} games, $N = 1{,}395{,}876$ scored asks
\cite{asking}.

\begin{table}[ht]
\centering
\caption{Ask calibration over $1{,}395{,}876$ scored asks, $300$ \usff{} games
\cite{asking}. ``Believed $p$'' is the number \code{pickAsk} decides on; ``observed'' is
the realised hit rate of those asks against the hands.}
\label{tab:calibration}
\small
\begin{tabular}{@{}lrrrr@{}}
\toprule
believed $p$ & $N$ & mean $p$ & observed & observed $-\ p$ \\
\midrule
0.0--0.1 & 126{,}303 & 0.007 & 0.0068 & $+0.0002$ \\
0.1--0.2 & 454{,}977 & 0.171 & 0.1681 & $-0.0028$ \\
0.2--0.3 & 717{,}054 & 0.221 & 0.2023 & $-0.0184$ \\
0.3--0.4 & 37{,}765  & 0.337 & 0.2912 & $-0.0461$ \\
0.4--0.5 & 24{,}124  & 0.440 & 0.4746 & $+0.0349$ \\
0.5--0.6 & 10{,}083  & 0.533 & 0.5205 & $-0.0130$ \\
0.6--0.7 & 8{,}181   & 0.625 & 0.6499 & $+0.0249$ \\
0.7--0.8 & 631       & 0.734 & 0.7242 & $-0.0095$ \\
0.8--0.9 & 78        & 0.810 & 0.7308 & $-0.0790$ \\
\textbf{0.9--1.0} & \textbf{16{,}680} & \textbf{1.000} & \textbf{1.0000}
  & $\mathbf{+0.0000}$ \\
\midrule
all & 1{,}395{,}876 & 0.2062 & 0.1953 & $-0.0109$ \\
\bottomrule
\end{tabular}
\end{table}

Two readings of Table~\ref{tab:calibration}, in order of load.

\textbf{The certainty tier is exact.} $16{,}680$ asks called certain hit $16{,}680$ times,
no exceptions. That is the property the whole policy is allowed to rest on --- row 9 keeps
the turn on a hit, so a certain hit is riskless material \emph{and} free tempo --- and it
holds at this scale without a single counterexample. It is also the strongest available
answer to the audit's question, because a false elimination that mattered would have to
show up here as a certain ask that missed.

\textbf{The overall bias is optimistic by about a point} ($-0.0109$: the engine's mean
belief exceeds the realised rate). That sign is the signature of an engine that is
\emph{not} over-inferring in the feared direction. An engine that wrongly concluded ``they
must have none'' would concentrate its remaining probability on the seats it had not
written off, and would therefore show hit rates \emph{above} its own claims in the low
buckets. Table~\ref{tab:calibration} shows the opposite in the three low buckets that
carry the mass --- $0.1$--$0.2$, $0.2$--$0.3$ and $0.3$--$0.4$, $1.21$ million asks
between them, all negative. The one low bucket that runs the other way is $0.0$--$0.1$, at
$+0.0002$: two ten-thousandths, on asks the engine already calls near-impossible, which is
not the concentration the feared error would produce. Among the buckets with mass, the two
most negative are $0.3$--$0.4$ ($-0.0461$) and $0.2$--$0.3$ ($-0.0184$), which is exactly
where the under-priced licence of Section~\ref{sec:underpricing} lives; the single most
negative row in the table is $0.8$--$0.9$ at $-0.0790$, but it rests on $78$ asks and
carries no weight.

The estimator here is a reliability decomposition of the familiar kind
\cite{murphy1973,dawid1982}: the engine's stated probability is the forecast, the realised
hit is the outcome, and the bucketed table is the reliability curve. Nothing is
recalibrated on the strength of it in this paper; Section~\ref{sec:substitution} explains
why the obvious recalibration is declined.

\section{What silence is actually worth}
\label{sec:silence}

Section~\ref{sec:absence} shows the engine does not use silence. This section asks what it
would be worth if it did. Two confounds are removed first: contained sets are excluded,
because the contained-book turn-pass deliberately manufactures repeat misses at seats that
provably hold nothing, and reuses the same card while doing it
\cite{containment,containedpaper}. Leaving them in would fabricate the effect the
measurement is looking for.

\begin{table}[ht]
\centering
\caption{Ask-level calibration by silence stratum, contained sets excluded
\cite{asking}. The third row is the owner's exact position.}
\label{tab:strata}
\small
\begin{tabular}{@{}p{0.44\textwidth}rrrr@{}}
\toprule
stratum & $N$ & believed & observed & error \\
\midrule
no prior miss, target \textbf{never asked} into the set & 1{,}176{,}652 & 0.2007 & 0.1927
  & $\mathbf{-0.0081}$ \\
no prior miss, target \textbf{has} asked into the set & 88{,}419 & 0.3473 & 0.4181
  & $\mathbf{+0.0708}$ \\
\textbf{1 miss by me, never asked back} (the owner's case) & 19{,}003 & 0.1399 & 0.0199
  & $\mathbf{-0.1200}$ \\
1 miss by me, target has asked into the set since & 21{,}726 & 0.3712 & 0.3792
  & $+0.0080$ \\
\bottomrule
\end{tabular}
\end{table}

\begin{table}[ht]
\centering
\caption{The set-level question directly: does the seat hold \emph{any} card of the set?
\cite{asking}}
\label{tab:setlevel}
\small
\begin{tabular}{@{}lrr@{}}
\toprule
position & $P(\text{holds} \ge 1 \text{ of the set})$ & $N$ \\
\midrule
no miss, silent & \textbf{0.663} & 263{,}016 \\
miss, then asked back & \textbf{0.761} & 6{,}180 \\
\textbf{miss, then silent} & \textbf{0.063} & 5{,}968 \\
two or more misses, then silent & \textbf{0.004} & 464 \\
\bottomrule
\end{tabular}
\end{table}

Three conclusions follow, and they are not the same conclusion.

\textbf{1. Silence on its own is genuinely near-uninformative.} Row 1 of
Table~\ref{tab:strata} is $0.8$ points of bias over $1.18$ million asks --- the largest
stratum in the study, and the one that corresponds to ``they simply have not asked into
this set''. The owner's principle is validated at scale: you cannot read anything into
``they didn't ask back'', taken alone.

\textbf{2. Silence \emph{after a miss} is a different object.} A seat you missed on that
then stays quiet holds a card of that set $6.3\%$ of the time (Table~\ref{tab:setlevel},
row 3), against $66.3\%$ for a silent seat you never missed on. Two or more misses then
silence takes it to $0.4\%$, on $464$ observations. The compound event is strongly
informative even though its parts are not, and the ask-level view agrees: the engine
believes $0.1399$ where the truth is $0.0199$. \emph{The engine under-uses this signal
rather than over-using it.}

\textbf{3. The engine is not making the mistake the owner was guarding against, in either
direction that matters.} Its one large error in this table is generosity. Note also row 2
of Table~\ref{tab:strata}, which is the mirror image and is the subject of
Section~\ref{sec:underpricing}: where a target \emph{has} published a licence, the engine
under-believes by $7.1$ points.

\subsection{Why the compound signal is not acted on anyway}
\label{sec:notacting}

Twelve points of ask-level bias and sixty points of set-level probability look like an
opportunity. The project declines it, and the reason is a measured one rather than a
temperamental one.

It is the same trap the companion document records \cite{concession}: a plausible small
signal read off the public log --- refuse to ask the opponent who would punish the
conceded turn --- cost up to $11.7$ points of win rate when acted on, because the row-6
confound made the information point the wrong way. Threat and opportunity are the same
fact under row 6, so the seat that is most dangerous to concede to is also the seat most
worth taking a card from. Silence is the \emph{weaker} of the two signals, and it is read
off the same channel.

If it is ever attempted, the discipline is fixed in advance and is the lab's own
\cite{botlab,v05paper}: duplicate deals, a byte-for-byte control that reproduces the
shipped policy exactly when the mechanism is switched off, and an information-free
perturbation of matched magnitude as the null. The last of those is the part that the
off-limits work established was necessary: without it, a mechanism that merely
\emph{changes} decisions can be mistaken for a mechanism that improves them.

\section{Two real gaps, and one pattern that looks like a third}
\label{sec:gaps}

Neither gap is the one the owner feared. Both are its opposite --- the engine under-reads
\emph{positive} evidence.

\subsection{A published licence is under-priced by eight points}
\label{sec:underpricing}

When the target has publicly asked into the set --- a live row-6 licence --- the engine
believes $0.2386$ and the truth is $0.3221$ ($N = 23{,}345$) \cite{asking}. It under-prices
its best asks by $8.4$ points.

The cause is the staleness of Section~\ref{sec:staleness}: the deal-time constraint that
records ``this seat holds a card of this set'' is dropped the moment it is satisfied or
exhausted, and \code{refinedHitProbability} iterates exactly those constraints.

Conditioning fixes it, and the conditioning is elementary. A live licence says the target
holds at least one unresolved card of the set, so under the engine's own per-card
estimates $q_j$ for the cards $j$ of set $B$ at target $t$,
\begin{equation}
\label{eq:condition}
\Pr\bigl(c \text{ at } t \;\big|\; \text{at least one of } B \text{ at } t\bigr)
\;=\;
\frac{q_c}{1 - \prod_{j \in B}\,(1 - q_j)} .
\end{equation}
Applied at strength $\lambda$ --- $\lambda = 0$ leaving the shipped probability untouched,
$\lambda = 1$ applying Equation~\eqref{eq:condition} in full --- the residual bias on the
licensed subset runs:

\begin{table}[ht]
\centering
\caption{The $\lambda$ ladder: residual calibration bias on the licensed subset
\cite{asking}.}
\label{tab:lambda}
\small
\begin{tabular}{@{}lrrrrrrr@{}}
\toprule
$\lambda$ & 0 & 0.25 & 0.40 & 0.50 & \textbf{0.60} & 0.75 & 1.00 \\
\midrule
bias & $-0.0835$ & $-0.0488$ & $-0.0280$ & $-0.0141$ & $\mathbf{-0.0002}$ & $+0.0206$
 & $+0.0554$ \\
\bottomrule
\end{tabular}
\end{table}

$\lambda = 0.60$ removes the bias almost exactly. That the correct strength is not $1$ is
itself informative: the licence is a fact about the deal, and by the time it is read the
seat may have been partly harvested, so the full conditioning over-corrects. The ladder is
monotone in $\lambda$ and crosses zero once, so the operating point is not a fitted
constant so much as a read-off. Section~\ref{sec:substitution} reports what
$\lambda = 0.60$ bought in play, which is the part that matters and is not what one would
guess.

\subsection{The conceded turn is priced with the wrong sign}
\label{sec:wrongsign}

The second gap belongs to the companion document \cite{concession} and is repeated here
because it belongs in the ask matrix. \code{missTarget}, \code{aimedTarget},
\code{valueContainedPass} and \code{signallingAsk} all rank the seat that receives the
turn by \textbf{hand size}, on the stated grounds that the surrendered turn is worth least
at the smallest hand. Over $12{,}376$ observed concessions, hand size correlates
$\mathbf{-0.147}$ with the cards that seat then actually takes --- negative in every
game-phase bucket --- against $\mathbf{+0.421}$ for its published reach, the sets it has
shown a row-6 basis in \cite{concession}.

The two gaps have the same root. Both are failures to use the row-6 licence: one on the
probability side, one on the concession side. The concession side has a fix that ships
(\code{defuse.ts}); the probability side has the fix derived in
Section~\ref{sec:underpricing} and not shipped, for the reason
Section~\ref{sec:substitution} measures.

\subsection{Why the same two or three sets get traded --- and why it is not a bug}
\label{sec:clustering}

The owner also reported watching games in which ``two or three half-suits get ask-traded
until it's declared, while the other half-suits never show up''. Both halves of that
report were tested, and they come apart \cite{asking}.

\textbf{The trading is real, and large.} After seat $B$ is asked for a card of set $H$,
misses, and takes the turn under row 10, $B$'s next ask goes into $H$ $\mathbf{88.1\%}$ of
the time when row 6 lets it ($N = 9{,}091$, $600$ games). The availability-matched control
--- every other legally-askable set at the same decisions --- is $\mathbf{21.4\%}$
($N = 203{,}723$). The provoking set beats the control $5$--$6\times$ in every bucket, so
this is not an artefact of what happens to be askable. Unconditionally the figure is
$36.0\%$, because $59\%$ of provocations land on a set $B$ holds no card of, where row 6
forbids the reply outright.

\textbf{But it is not a reciprocation rule. It is the certainty rule.}

\begin{table}[ht]
\centering
\caption{What the reply into the provoking set actually is \cite{asking}.}
\label{tab:reply}
\small
\begin{tabular}{@{}p{0.62\textwidth}r@{}}
\toprule
of $B$'s replies into the provoking set $H$ & share \\
\midrule
replies into $H$ that are the argmax on hit probability & \textbf{99.1\%} \\
replies into $H$ that are \textbf{certain hits}, after a miss & 60.2\% \\
\quad \dots{} after a hit & 94.3\% \\
$B$ would have chosen $H$ \emph{anyway}, before being provoked & \textbf{70.1\%} \\
decisions the provocation actually switched into $H$ & \textbf{18.1\%} \\
\bottomrule
\end{tabular}
\end{table}

An ask into $H$ publishes that the asker holds a card of $H$ (row 6) and lacks the named
one (row 7). Combined with what $B$ already knows, that very often \emph{locates a card
exactly} --- and \code{certaintyBonus} makes a certain hit dominate the sort, while row 9
keeps the turn. \textbf{The bot is not asking back; it is taking a card it can see, and
the ask it was just asked is what made it visible.} It could not be doing anything else:
the tradition bans prearranged conventions \cite{develin}, and nothing in this engine
decodes an incoming ask as a signal --- only as the two facts rows 6 and 7 make it
publish.

Ablations confirm the ordering, and they are decision-level rather than
outcome-level --- the unit the project holds to be the honest one for a question of this
shape \cite{inertpaper}. The $88.1\%$ falls to $84.4\%$ with \mbox{\code{defuse: 0}}, to $72.3\%$
with refined inference off, and to $72.0\%$ with constraints off --- still $3.4\times$ the
control. The defusal term contributes about four points of it and is not the cause; the
\code{-nodefuse} arms reproduce committed HEAD byte-for-byte on the ask path, which is
what licenses that conclusion \cite{asking}.

\textbf{The ``other sets never show up'' half is false at game scale and true locally.}
Per game the top three sets take $49.2\%$ of asks against a uniform $33.3\%$, an HHI of
$0.1358$ against $0.1111$ --- real but modest --- and only $0.55$ of $9$ sets receive no
ask at all. But in any window of $10$ consecutive asks only $\mathbf{2.96}$ distinct sets
appear, against $\mathbf{5.98}$ for the same games' asks randomly permuted, and $75.3\%$
of windows show three or fewer sets against $0.4\%$ permuted. The owner is reporting a
local view accurately: the clustering is \emph{temporal}, not a coverage failure. Every
set does get its turn; they just take turns in long runs.

\textbf{Does the clustering cost anything? Not identifiably.} The raw correlation between
a game's ask concentration and its set outcome is $-0.443$, but ask \emph{count}
correlates $+0.929$ with sets, and controlling for it collapses the partial correlation to
$-0.074$ ($-0.008$ on the roster arm). Ask count is plausibly a mediator as well as a
confound --- games that run long have more asks \emph{and} more sets --- so this
identifies nothing causal in either direction, and no intervention is warranted on this
evidence. The honest statement is that the observed behaviour is explained, not that it is
optimal.

\section{The holistic ask matrix}
\label{sec:matrix}

The owner asked for the whole reasoning matrix for choosing an ask, not only the silence
question. Table~\ref{tab:matrix} is it: twelve factors, each with the rule it rests on,
the sign it should carry, and whether the engine prices it today \cite{asking}.

\begin{table}[p]
\centering
\caption{The ask matrix \cite{asking}. ``Priced today?'' is a statement about the shipped
engine, and the two entries that are not ``yes'' are the two gaps of
Section~\ref{sec:gaps}.}
\label{tab:matrix}
\footnotesize
\begin{tabular}{@{}c p{0.19\textwidth} p{0.25\textwidth} p{0.13\textwidth} p{0.27\textwidth}@{}}
\toprule
\# & factor & grounding & sign & priced today? \\
\midrule
1 & \textbf{hit probability} & rows 9/10 --- a hit takes a card and keeps the turn; a miss
  ends it & \textbf{+}, dominant & \textbf{yes, well.} \code{wHit}$\cdot p$ is 70 of 100 of
  the score; overall bias $-1.1$ pt; the certainty tier is exact \\
2 & \textbf{certainty vs.\ probability} & a certain hit is riskless material \emph{and}
  free tempo & \textbf{+}, must dominate & \textbf{yes.} \code{certaintyBonus} $\ge 20$
  required of every style; the contained pass refuses to displace one \\
3 & \textbf{target's published licence in the asked set} & row 6 is the only ask licence,
  so an ask publishes ``I hold one of these'', permanently & \textbf{+} on both the
  probability and the value of taking the card & \textbf{partly --- the biggest gap.} The
  \emph{choice} is priced (\code{defusalBonus}: $+1.50$/$+1.65$ sets for the prototype
  term, $+1.43$/$+1.58$ for the shipped implementation). The
  \emph{probability} is not: \S\ref{sec:underpricing}, $-8.4$ points \\
4 & \textbf{what my ask publishes about me} & rows 6 and 7, both permanent
  & \textbf{--}, small & \textbf{yes, as a tie-break only} (\code{leakEpsilon}).
  Deliberately narrow: the best ask into a strong set is usually the ask that completes
  it \\
5 & \textbf{who receives the turn on a miss} & row 10 & \textbf{--} on the miss branch ---
  but see 6 & \textbf{wrong sign.} \S\ref{sec:wrongsign} \\
6 & \textbf{threat and opportunity are the same fact} & row 6 makes it structural; threat
  correlates $+0.249$ with the best available $p$, whose mean is $0.471$ at high-threat
  seats against $0.326$ elsewhere & \textbf{the naive veto is negative}
  & \textbf{yes, correctly.} Avoidance costs up to $11.7$ points; \code{defuse.ts} inverts
  it \\
7 & \textbf{defusal credit} & a hit strips the licence and row 9 keeps the turn
  & \textbf{+} & \textbf{yes}, \mbox{\code{defuse: 1}} across the roster, replicated in a foreign
  engine \cite{crossplay} \\
8 & \textbf{is the set contained?} & no opponent can ask into a one-team set, and declaring
  it hands it back & \textbf{flips the calculus}: the ask stops being about cards and
  becomes purely about aiming the turn & \textbf{yes, own branch} --- but its
  \code{gain} is built on factor 5's backwards proxy \\
9 & \textbf{progress toward the set} & row 12 needs all six holders named
  & \textbf{+}, modest & \textbf{yes}, \code{wProgress} \\
10 & \textbf{narrowing value of a miss} & row 17 makes every ask public; a miss on a
  two-candidate card pins it & \textbf{+}, small & \textbf{yes}, \code{wNarrow} --- the
  only term that pays for a miss \\
11 & \textbf{clinch state} & the game ends at 5 sets; a wrong declare gifts one
  & \textbf{asymmetric, large late} & \textbf{declares yes, asks no.} Nothing in
  \code{pickAsk} reads the score. Defensible: an ask cannot lose a set outright \\
12 & \textbf{silence} & not a rule --- an absence & \textbf{--}, weak alone
  (\S\ref{sec:silence}) & \textbf{ignored, and that is correct} \\
\bottomrule
\end{tabular}
\end{table}

Summarised: \textbf{priced well}, factors 1, 2, 4, 7, 9, 10, 11; \textbf{wrong sign},
factor 5, with factor 8 inheriting it; \textbf{under-priced}, factor 3; \textbf{correctly
ignored}, factor 12. The matrix's own headline is that the engine's errors are
concentrated in one place --- what a published row-6 licence is worth --- and that they
all run in the direction of under-use.

\subsection{The one thing this matrix must not say}
\label{sec:offlimits}

The owner pairs ``whom to ask'' with ``keeping in mind the off-limits rule'': the idea
that some opponents should be off limits because the turn you concede on a miss would be
punished. This project's measured position is that the off-limits \emph{prescription} ---
refuse to ask the dangerous seat --- \textbf{loses $4.5$ to $11.7$ points of win rate},
with the paired set-difference falling monotonically with appetite \cite{concession}. Row
6 is the reason: an opponent is dangerous \emph{because} it has published a basis, and a
published basis is also the strongest available evidence that an ask into that set will
hit. Threat and opportunity are the same fact.

The threat \emph{model} is kept and the prescription inverted --- ask them, to strip the
card their reach rests on --- which is \code{defuse.ts}. Two measurements of it are on
record and both belong here: the prototype term, hooked into a mirror of \code{pickAsk},
is worth $+1.50$ and $+1.65$ sets per duplicate pair on two $600$-pair held-out banks,
while the \emph{shipped} implementation re-measured on those same two banks at $400$ pairs
lands lower, at $+1.43 \pm 0.32$ and $+1.58 \pm 0.32$; the gap sits inside the intervals
and is what a smaller bank plus a real integration should look like. The gain is
replicated in a foreign engine \cite{concession,crossplay}. The only asks for which the
original avoidance rule is right
are the ones that \emph{cannot hit}, and those are exactly the asks where the engine still
ranks the receiving seat by hand size (Section~\ref{sec:wrongsign}).

\section{Two corrections for the same evidence, and what stacking them costs}
\label{sec:substitution}

Section~\ref{sec:underpricing} removes a real $8.4$-point bias with a principled,
knob-free conditioning. The obvious expectation is that a better probability plays better.
It does --- \emph{but only if the engine is not already being paid for the same evidence.}

\subsection{The design, and why the common baseline is the whole point}
\label{sec:commonbaseline}

All arms mirror \code{pickAsk} as it now stands, defusal term included, with the zero arm
verified to reproduce the shipped policy byte-for-byte --- the lab's standard control
discipline \cite{botlab,v05paper}. The measurement is $400$ duplicate pairs, one bank,
and --- the design decision that carries the result --- \textbf{one common baseline}
(\mbox{\code{defuse: 0}}, $\lambda = 0$), so that the three treatments are directly comparable to
each other and not merely each to itself \cite{asking}.

\begin{table}[ht]
\centering
\caption{The three-way substitution, $400$ duplicate pairs, one bank, one common baseline
(\mbox{\code{defuse: 0}}, $\lambda = 0$) \cite{asking}. Intervals are $\pm 1.96$ SE on the paired
set-difference. See Section~\ref{sec:scale} on the $400$-pair scale.}
\label{tab:substitution}
\small
\begin{tabular}{@{}lrr@{}}
\toprule
arm & win rate & paired set-difference \\
\midrule
baseline --- neither & 50.00\% & $0.0000 \pm 0.0000$ \\
\textbf{defusal alone} (\mbox{\code{defuse: 1}}, $\lambda = 0$) & 64.00\%
  & $\mathbf{+1.915 \pm 0.315}$ \\
\textbf{calibration alone} (\mbox{\code{defuse: 0}}, $\lambda = 0.6$) & 65.25\%
  & $\mathbf{+2.068 \pm 0.293}$ \\
calibration alone, over-corrected ($\lambda = 1.0$) & 61.63\% & $+1.620 \pm 0.327$ \\
\textbf{both together} (\mbox{\code{defuse: 1}}, $\lambda = 0.6$) & 60.63\%
  & $\mathbf{+1.565 \pm 0.317}$ \\
\bottomrule
\end{tabular}
\end{table}

\subsection{The result}
\label{sec:subresult}

Three things follow from Table~\ref{tab:substitution}, in order of how much they change
what one would do.

\textbf{1. The two mechanisms look like substitutes rather than complements, and the
ordering is suggestive rather than established.} Either alone is worth about two sets per
duplicate pair, and the stack measures below either alone. Taken at face value that says
adding the principled fix on top of the shipped heuristic would have made the engine
worse.

The claim has to be stated at exactly the strength the design supports, which is less than
the face value. \textbf{Every interval in Table~\ref{tab:substitution} is on an arm
against the common baseline; no interval is published on any \emph{difference} between two
arms}, and the differences are what the substitution claim is about. Nor can the arm
intervals be differenced by eye: the three treatments run on one bank against one
baseline, so their errors are positively correlated and a properly paired arm-to-arm
interval would be \emph{narrower} than any naive combination of the two published
half-widths. That is the design's strength --- a common bank and a common baseline make
this a paired comparison, far sharper than measuring each treatment on its own bank and
subtracting --- but the paired interval was not computed, so the gap is unquantified in
either direction.

What can be said about magnitude is discouraging at this $N$. The larger gap, calibration
alone against the stack, is $0.503$ sets; the smaller, defusal alone against the stack, is
$0.350$. Section~\ref{sec:scale} puts each arm's own $80\%$-power minimum detectable
effect at $0.42$--$0.47$ sets, so the smaller gap sits below the floor these arms were
measured at and the larger one barely clears it --- and the empirical floor at $N = 400$
is itself only pinned to $0.36$--$0.50$. The pairing buys back some of that, and the
ordering is consistent across both comparisons and reinforced by the $\lambda = 1.0$ row
below, which is why the result is worth reporting. But \textbf{it is a suggestive ordering
at this bank size, not an established one}, and the honest next step is to re-run the same
five arms with the paired arm-to-arm differences and their intervals reported, at a bank
where $0.35$ sets is comfortably resolvable.

What the design \emph{does} establish is that this question is askable at all. Only a
common-baseline design can see substitution: measured each against its own baseline --- the
natural thing to do, and the thing an incremental development process does by default ---
the calibration fix looks like a flat $+0.04 \pm 0.41$ rather than like something competing
with a term already shipped.

\textbf{2. The mechanism is one reordering reached two ways.} Equation~\eqref{eq:condition}
multiplies every card of one set at one seat by the same factor, because the denominator
depends on the set and the seat and not on the card. It therefore \emph{cannot} reorder
asks within a (set, seat) pair; all it can do is promote licensed asks against unlicensed
ones. That is exactly what \code{defuse} does, by a different route --- a bonus on the
seats whose published reach makes them worth stripping. Two terms that promote the same
asks are not additive, and the $\lambda = 1.0$ row is the same shape seen from inside a
single mechanism: over-shoot the promotion alone and the gain falls from $+2.068$ to
$+1.620$, which is what stacking also produces.

\textbf{3. The calibration error is real and worth fixing eventually --- just not here.}
\code{refinedHitProbability} is also read by the declare planner, by
\code{valueContainedPass}, and by any future search or exploitability sweep, and in none of
those is a defusal bonus quietly compensating for it. It is a bug fix rather than an
appetite; it needs no style knob; it would help every consumer of the probability rather
than only the ask ranker.

\textbf{The recommendation is to ship neither change now.} \mbox{\code{defuse: 1}} stays, on the
strength of its gauntlet against opponents that also defuse and its foreign-engine
replication \cite{concession,crossplay}, neither of which the calibration fix has. The
calibration fix is recorded with its $\lambda$ ladder as the better long-term shape.
Swapping the two is a deliberate experiment for its own day, not a change to make in
passing.

\subsection{The scale caveat, and a refuted explanation}
\label{sec:scale}

Table~\ref{tab:substitution}'s numbers come from a $400$-pair bank, and that is this
instrument's thin end. The project has since characterised the instrument: for arms that
differ at every decision the per-pair standard deviation is $3.15$--$3.44$ sets, which
puts the $80\%$-power minimum detectable effect at about $0.47$ sets at $N = 400$ and
about $0.33$ at $N = 800$ \cite{concession}.

That generic figure must not be applied to this cell by eye, and the companion document
names the substitution as a trap it fell into itself: the minimum detectable effect
depends on \emph{that cell's own} standard deviation, which varies across the repository
by up to a half \cite{concession}. So the floor here is computed from
Table~\ref{tab:substitution}'s own published half-widths. A half-width of $h$ at
$N = 400$ implies a per-pair SD of $h\sqrt{N}/1.96$, giving $3.21$ for the defusal arm,
$2.99$ for the calibration arm, $3.23$ for the stack and $3.34$ for the over-corrected
arm; the corresponding $80\%$-power minimum detectable effects are $0.45$, $0.42$, $0.45$
and $0.47$ sets. Each arm clears its own floor by three to five times, so each
arm's effect against the common baseline is resolved. What is \emph{not} resolved is the
spacing between the arms: see Section~\ref{sec:subresult}.

Two properties of the instrument at $N = 400$ failed to replicate, and both belong here.
The coverage is unsettled --- an independent replication found $10.8\%$ false positives
against a nominal $5\%$, pooling to $8.0\%$ with a confidence interval that excludes
$5\%$, and whether that is real under-coverage or noise is recorded as unresolved. So is
the empirical floor: a planted effect of $0.393$ sets was detected in $84.7\%$ of banks on
the first run and $75.0\%$ on the replication, pooling to $81.7\%$, which leaves the
$N = 400$ floor somewhere in $0.36$--$0.50$ and not pinned tighter by that design ---
the $N = 800$ floor does replicate \cite{concession}. Intervals at this scale should be
read with both in mind.

More directly: \textbf{the $+1.92$-scale figures on this bank are high.} The companion
document re-ran the defusal quantity at $4{,}000$ duplicate pairs through a single code
path, and the same edge measures $+1.6225 \pm 0.1016$, replicated independently at
$+1.6595 \pm 0.1022$, against Table~\ref{tab:substitution}'s $+1.915 \pm 0.315$
\cite{concession}. The gap is ordinary sampling variation, and the record of how that was
established is worth keeping, because the first answer offered was wrong: it was proposed
that this paper's mirror harness lacks \code{pickAsk}'s certain-hit gate and therefore
measured a different quantity. That explanation was \emph{tested}, by running the ungated
harness, and \textbf{refuted} --- $+1.6850 \pm 0.3358$, indistinguishable from the gated
value \cite{concession}. The harness difference is not the explanation; the bank size is.

None of this disturbs the ordering Table~\ref{tab:substitution} exists to establish,
because all five arms share one bank, one baseline and one code path, and the substitution
result is an ordering among them rather than an absolute level. It does mean the
\emph{levels} in that table should be read as this bank's, and the $4{,}000$-pair figures
preferred wherever the same quantity has been measured there.

\subsection{The general lesson, stated once}
\label{sec:lesson}

\emph{A term that corrects a probability and a term that rewards the same evidence are not
additive.} Measuring the second after the first has shipped will show nothing --- the
$+0.04 \pm 0.41$ of Section~\ref{sec:subresult} is exactly what that looks like --- and
stacking them can subtract. Any future mechanism that promotes asks at seats with a
published row-6 licence is competing with \code{defuse} for the same effect, and must be
measured against a baseline with \code{defuse} \emph{off} before it can be believed.

The lesson generalises past this engine. Wherever a heuristic has been tuned into a
system and a principled correction for the same underlying error is proposed later, the
default comparison --- new thing versus current system --- is the one comparison that
cannot detect substitution. The common baseline costs one extra arm and is the only design
that answers the question actually being asked.

\section{Threats to validity}
\label{sec:threats}

\begin{enumerate}
\item \textbf{The soundness audit is a code audit, not a proof.} Table~\ref{tab:audit}
  enumerates the elimination paths present in the implementation and licenses each against
  a rule; it cannot establish that no future path will be added, and it says nothing about
  paths in other engines. Its strongest empirical backing is the exact certainty tier of
  Table~\ref{tab:calibration}: a false elimination that mattered would have to surface
  there.

\item \textbf{Calibration is measured against this roster's own play.} The
  $1{,}395{,}876$ asks of Table~\ref{tab:calibration} come from $300$ self-play \usff{}
  games. The distribution of positions is the roster's, so the reliability curve is
  conditional on the population that generated it. A different opponent population would
  produce different bucket masses and could produce different biases; nothing here
  measures that.

\item \textbf{The silence strata are observational.} Table~\ref{tab:strata} conditions on
  events the engine's own policy produced --- which seats it chose to ask, and when. The
  strata are therefore not randomised, and the $6.3\%$ figure of
  Table~\ref{tab:setlevel} describes seats the engine chose to miss on rather than seats
  in general. Contained sets are excluded for exactly this reason
  (Section~\ref{sec:silence}); other selection effects of the same shape are possible and
  unmeasured, which is one more reason Section~\ref{sec:notacting} declines to act on the
  signal without a duplicate-deal experiment.

\item \textbf{$\lambda = 0.60$ is read off one ladder.} Table~\ref{tab:lambda} is a
  post-hoc sweep on the licensed subset of the same games the bias was measured on. The
  ladder is monotone and crosses zero once, so the operating point is not fragile in
  shape, but it has not been validated out of sample and is not offered as a constant to
  ship.

\item \textbf{Table~\ref{tab:substitution} is a $400$-pair bank.}
  Section~\ref{sec:scale} states this in full: each arm clears its own detection floor
  against the common baseline, the levels are high against a $4{,}000$-pair
  re-measurement of the one quantity that has one, and both the interval coverage and the
  empirical floor at this $N$ failed to replicate. The paper does not claim the levels,
  and it claims the ordering only as suggestive: no interval is published on any
  difference between arms, and the smaller of the two gaps the substitution reading rests
  on is below these arms' own floor (Section~\ref{sec:subresult}).

\item \textbf{The substitution result is one operating point.} One $\lambda$, one defusal
  appetite, one roster, one bank. That the two mechanisms appear to substitute at
  $(\lambda, \code{defuse}) = (0.6, 1)$ does not establish the shape of the interaction
  surface elsewhere on it, and no arm here measures a reduced defusal appetite alongside
  the calibration fix --- which is the obvious next experiment and is not run.

\item \textbf{Ask clustering is explained, not priced.} Section~\ref{sec:clustering}
  identifies the mechanism behind the $88.1\%$ and shows the concentration is temporal.
  The partial correlation of $-0.074$ is explicitly not a causal estimate --- ask count is
  plausibly a mediator as well as a confound --- so nothing here says whether the
  clustering helps, hurts, or neither.

\item \textbf{Self-play, one engine, one rule set.} The usual transfer caveats apply in
  full \cite{v05paper,v10paper}. The one external check available is the foreign-engine
  replication of defusal \cite{crossplay}, which covers the mechanism of factor 7 and none
  of the inference results in this paper.
\end{enumerate}

\section{Conclusion}
\label{sec:conclusion}

The owner's worry was that a strong inference engine would slide from ``no evidence they
hold one'' to ``evidence they hold none''. The engine does not, and the reason is
structural: absence is not an event, and the layer that would have to act on one is a
switch over events that happened. The measured residual runs the other way --- $14.0\%$
believed against $2.0\%$ true in exactly the position the worry describes --- and the
calibration table says the same thing over $1{,}395{,}876$ asks: the engine is optimistic
by about a point, which is the wrong sign for the feared error and the right sign for an
engine that under-uses what it knows. The certainty tier, the one place where the policy
is allowed to rest its whole weight on the inference, is exact over $16{,}680$ asks.

What the study found instead is that the engine's real errors are all one error seen from
different angles: it under-uses the row-6 licence. It under-prices the probability of a
hit at a seat that has published a basis by $8.4$ points, because the constraint recording
that basis is dropped the moment it is confirmed. It ranks the seat receiving a conceded
turn by hand size, a proxy that correlates $-0.147$ with what that seat then takes against
$+0.421$ for the licence. And the behaviour that looked most like a bug from the
outside --- two or three sets traded back and forth --- turns out to be the engine using
the licence \emph{correctly}: $99.1\%$ of those replies are the hit-probability argmax, and
the bot is not reciprocating but collecting certainties that the incoming ask made
visible.

The result that changed what the project would do is the last one. A principled
correction for the under-priced licence is worth $+2.068 \pm 0.293$ sets per duplicate
pair on a common baseline; the shipped heuristic that rewards the same evidence is worth
$+1.915 \pm 0.315$; both together are worth $+1.565 \pm 0.317$. Two fixes for one error
measure lower stacked than singly --- an ordering that is suggestive at this bank and not
established, because no interval was published on any difference between the arms and the
smaller gap is below their own detection floor. The design claim is the firm one: the only
comparison that can \emph{see} substitution is one that measures both mechanisms against a
baseline carrying neither. The recommendation is to ship neither change now, and to record the
calibration fix --- with its ladder, its bank size, and the refuted explanation of its
scale --- as the better long-term shape for whoever swaps them deliberately.

\medskip
\noindent\textbf{Reproducibility.} The engine is deterministic; every number in this paper
is a field of a committed probe result in \cite{asking} or \cite{concession}. The
calibration of Section~\ref{sec:calibration} is \code{scripts/probe-inference.mjs}; the
strata of Section~\ref{sec:silence} are \code{scripts/probe-inference-strata.mjs}; the
$\lambda$ ladder of Section~\ref{sec:underpricing} is \code{scripts/probe-licence.mjs};
the clustering of Section~\ref{sec:clustering} is \code{scripts/probe-concentration.mjs}
with its window and outcome companions; and Table~\ref{tab:substitution} is
\code{scripts/probe-licence3.mjs}, whose \code{0 lic3 0 0} invocation is the byte-exact
control that must print $0.0000 \pm 0.0000$ \cite{asking,repo}.

\begin{thebibliography}{99}
\footnotesize

\bibitem{murphy1973}
A.~H.~Murphy.
\newblock A new vector partition of the probability score.
\newblock \emph{Journal of Applied Meteorology}, 12(4):595--600, 1973.
The reliability/resolution decomposition Table~\ref{tab:calibration} is an instance of.

\bibitem{dawid1982}
A.~P.~Dawid.
\newblock The well-calibrated Bayesian.
\newblock \emph{Journal of the American Statistical Association}, 77(379):605--610, 1982.

\bibitem{develin}
M.~Develin.
\newblock \emph{The Ten Best Card Games You've Never Heard Of}, ch.~9 (Canadian Fish).
\newblock \url{http://www.bantha.org/~develin/cardgames.html}.
Independent attestation of the absorbing property of a contained set, and of the tradition's
ban on prearranged conventions --- which is why nothing in this engine decodes an ask as a
signal.

\bibitem{pagat}
J.~McLeod (ed.).
\newblock \emph{Literature (Canadian Fish)}.
\newblock Pagat card-game rules archive,
\url{https://www.pagat.com/quartet/literature.html}.

\bibitem{v05paper}
A.~Hsieh.
\newblock FishAI v0.5: Measuring play style without skill in a 54-card Literature
variant.
\newblock FishAI project; companion paper, 2026. The roster, the duplicate-deal protocol,
and the byte-for-byte control discipline every measurement here inherits.

\bibitem{v10paper}
A.~Hsieh.
\newblock FishAI v1.0: When best-response adaptation is worth less than nothing.
\newblock FishAI project; companion paper, 2026.

\bibitem{v15paper}
A.~Hsieh.
\newblock FishAI v1.5: A bit-budget memory ladder as an honest difficulty axis.
\newblock FishAI project; companion paper, 2026. The bounded fact pool whose
\code{lacks-card} and \code{no-basis} facts Section~\ref{sec:audit} audits.

\bibitem{containedpaper}
A.~Hsieh.
\newblock The contained book: an absorbing resource measured to be worth nothing.
\newblock FishAI project; companion paper, 2026. The turn-pass whose repeat misses
Section~\ref{sec:silence} excludes as a confound.

\bibitem{inertpaper}
A.~Hsieh.
\newblock The inert axis: when a style parameter is wired, swept, and never reached.
\newblock FishAI project; companion paper, 2026. The argument that decision-level
divergence, not outcome, is the honest unit for a question of this shape.

\bibitem{obspaper}
A.~Hsieh.
\newblock What a public log reveals: observing play style in Literature (Canadian Fish).
\newblock FishAI project; companion paper, 2026. The specification of the public channel
this paper's inferences are drawn from.

\bibitem{asking}
A.~Hsieh.
\newblock ASKING.md --- who to ask, for what, and what silence is worth.
\newblock FishAI project internal report, August 2026. \S1 the elimination-path audit,
\S2 the calibration over 1{,}395{,}876 asks, \S3 the silence strata, \S4.1 the licence
under-pricing and its $\lambda$ ladder, \S4.3 ask clustering, \S5 the ask matrix, \S6 the
three-way substitution.

\bibitem{concession}
A.~Hsieh.
\newblock CONCESSION.md --- FishAI v2.0: the three-sided ask.
\newblock FishAI project internal report, August 2026. \S1.2 the two concession
estimators, \S2 off-limits refuted, \S3 defusal, \S8.1 the triangle refuted and the
$400$-pair/$4{,}000$-pair correction, \S8a the detection floor.

\bibitem{containment}
A.~Hsieh.
\newblock CONTAINMENT.md --- the contained-book result.
\newblock FishAI project internal report, August 2026.

\bibitem{crossplay}
A.~Hsieh.
\newblock CROSSPLAY.md --- cross-engine play against a foreign Literature implementation.
\newblock FishAI project internal report, August 2026. \S3, the foreign-engine replication
of the defusal gain.

\bibitem{rules54}
FishAI project.
\newblock RULES\_US54.md --- the ``US student'' 54-card variant.
\newblock Repository rules specification with derived consequences and test vectors, 2026.
Decision-table rows cited by number throughout.

\bibitem{botlab}
FishAI project.
\newblock BOT\_LAB.md --- play-style spectrum: research, experimental design, and metrics.
\newblock Repository document; \S5 experimental design, including duplicate-deal pairing
and the null-control requirement, 2026.

\bibitem{repo}
FishAI project.
\newblock FishAI repository: engine, bots, simulation laboratory, and results site.
\newblock Inference layer \code{lib/engine/bots/knowledge.ts} (the ingestion \code{switch}
and \code{refinedHitProbability}); the threat model
\code{lib/engine/bots/threat.ts}, whose header records the constraint-staleness defect of
\S\ref{sec:staleness}; the defusal term \code{lib/engine/bots/defuse.ts}; the ask ranker
\code{lib/engine/bots/decide.ts}; probes under \code{scripts/}, 2026.

\end{thebibliography}

\end{document}
